\scanpage{69}
\mspage{22}
\uline{Gluing theorem for tubes} (\unsure{1\textsuperscript{st}|1\textsuperscript{st}|firs\textsuperscript{t}} version) Consider the diagram of type $\widetilde I$, formed by the $X_i^*,V_{i,j}^*,V_{i,j}$ ($i,j\in I$, $j$ a successor of $i$) and the three types of arrows appearing \unsure{there} (in solid lines) in (17). Then $X$ identifies with the \deleted{topos} \inserted{topos} $\underset{\longrightarrow}{\mathrm{Lim}}$ of these toposes.

\margin{For this statement to be able to take on a precise meaning, one \unsure{must} probably assume $X$ \uline{sober}, i.e. determined by $\operatorname{Top}(X)\ldots$}

\deleted{Lacking proper definitions for tubular neighborhoods, I am going to give a proof, when one \deleted{\unsure{sets}} $X$ replaces \deleted{\unsure{the|each}} $V_{ij}$ by \inserted{genuine} genuine open subsets of $X_j$\unsure{,} of the form $\bar V_{i,j}\setminus R_{i,j}\cap\bar V_{i,j}$, where $\bar V_{i,j}$ is a neighborhood of $X_i$ in $X_j$, and one sets $\unsure{U_{ij}^*|V_{ij}^*}=U_{ij}\setminus X_i^*$. \deleted{Then} \deleted{some \unsure{neighborhoods} $V_{ij}$}}

We are going to give a heuristic proof of it below. But before doing so, I am going to generalize it, into a statement which makes it possible to express \unsure{als\textsuperscript{o}|als\textsuperscript{o}} the ``canonical toposes'' associated with the \unsure{stratificatio\textsuperscript{n}|stratificatio\textsuperscript{n}}, in terms of the inductive system of type $\widetilde I$ formed by the elementary toposes.

\margin{68}
